\chapter{Hardware Transactional Memory – Design and Simulation}
\index{hardware transactional memory!design}\index{hardware transactional memory!simulation}\index{RTM}\index{TME}\index{cache coherence}\index{capacity aborts}\index{fallback paths}

Hardware transactional memory, or HTM, is the most direct attempt to make transactions disappear into the memory hierarchy. Instead of maintaining read sets, write sets, undo logs, and validation logic in software, the processor extends speculation and coherence so that a region of ordinary loads and stores either commits atomically or aborts and restores the pre-transaction state.

This chapter studies HTM as an architectural idea and as a practical programming interface. The running example remains the two-account bank transfer: a transfer must preserve the account floor and the money-conservation invariant,

$$
\operatorname{total}(b)=b_0+b_1=\operatorname{constant}.
$$

The central lesson is deliberately modest: HTM can make the fast path extremely cheap, but it is a best-effort mechanism. Correctness lives in the fallback path.

\section{Extending Cache Coherence for TM}

Herlihy and Moss proposed the basic hardware idea: add transactional states to a MESI-like cache-coherence protocol. A transactional load marks a cache line as part of the transaction's read set. A transactional store marks a line as part of its write set and delays publication until commit. Their proposal used states such as XCOMMIT and XABORT to distinguish speculative lines from ordinary modified or shared lines.

The cache hierarchy then becomes the transaction manager. It tracks which lines have been read, which lines have been written, and whether another core sends a coherence message that invalidates a transactional assumption. A conflicting invalidation to a transactional line is an abort signal. The hardware discards the speculative writes and restarts execution at the checkpointed program counter.

The expensive requirement is undo. If the processor has overwritten a cache line speculatively, it must still be able to restore the pre-transaction value. A simple design keeps both the new value and enough old state inside the cache. This is fast when the transaction fits in the tracked cache structures and brittle when it does not. Evicting a transactional line may lose either the value needed for rollback or the metadata needed for conflict detection. The conservative answer is a capacity abort.

This explains why coherence-based HTM tends to work best for short, cache-resident critical sections. It is ideal for the common case in which a lock protects a small update, such as the bank transfer touching two account records. It is much less robust for long scans, large read sets, or transactions that call into the operating system.

\begin{figure}[htbp]
\centering
\begin{tikzpicture}[
    box/.style={draw, rounded corners, align=center, minimum width=2.8cm, minimum height=0.9cm},
    smallbox/.style={draw, rounded corners, align=center, minimum width=2.3cm, minimum height=0.8cm},
    arr/.style={-{Latex}, thick}
]
\node[box] (begin) {XBEGIN\\checkpoint PC};
\node[box, right=1.4cm of begin] (track) {L1 tracks\\read/write lines};
\node[box, right=1.4cm of track] (body) {speculative\\loads/stores};
\node[box, right=1.4cm of body] (commit) {XEND\\publish writes};

\node[smallbox, below=1.1cm of track] (inv) {coherence\\invalidation};
\node[smallbox, below=1.1cm of body] (evict) {capacity\\eviction};
\node[box, below=1.1cm of $(inv)!0.5!(evict)$] (abort) {abort\\restore checkpoint};

\draw[arr] (begin) -- (track);
\draw[arr] (track) -- (body);
\draw[arr] (body) -- (commit);
\draw[arr] (inv) -- (track);
\draw[arr] (track) -- node[left, align=center]{conflict} (abort);
\draw[arr] (evict) -- node[right, align=center]{lost tracking} (abort);
\draw[arr] (body) -- (evict);
\draw[arr] (abort.west) -| (begin.south);
\end{tikzpicture}
\caption{Coherence-based HTM treats invalidations and capacity loss as abort signals. Commit is cheap only while the read and write sets remain tracked in cache.}
\label{fig:htm-coherence-abort}
\end{figure}

The figure is also a useful mental model for software design. Every extra cache line touched inside the transaction becomes part of the speculative footprint. Debug counters, logging buffers, allocator metadata, and accidental reads of shared state all increase the chance of abort. In HTM programming, minimizing the hardware footprint is not merely an optimization; it is part of making progress likely.

\section{Modern HTM Implementations}

Commodity HTM is not one mechanism. It is a family of best-effort mechanisms with similar semantics and different interfaces. Intel Restricted Transactional Memory, or RTM, exposes instructions such as \texttt{xbegin}, \texttt{xend}, and \texttt{xabort}. ARM Transactional Memory Extension, or TME, exposes similar architectural concepts with implementation-defined tracking and diagnostics. IBM's Blue Gene/Q and POWER implementations are the longevity case study: they show that coherence-based HTM can be made practical when the architecture and platform expose a sufficiently mature interface.

The shared semantic contract is simple: a hardware transaction may commit, but software must not depend on eventual commit. It may abort because another core conflicts with its read or write set. It may abort because the transaction exceeds implementation capacity. It may abort because it executes an unsupported instruction, takes an interrupt, faults on a page, or hits a microarchitectural limit that is not visible in the source code.

This best-effort rule distinguishes HTM from the software backends studied earlier. In NOrec, TL2, TinySTM, SwissTM, JVSTM, MVLog, and the persistent and distributed designs, the implementation owns the retry protocol. With HTM, the processor owns only the fast attempt. The program must provide a correct fallback path.

The simplest fallback is a global lock: try HTM a few times, then serialize. The companion codebase includes TSX-SGL as the natural instance of this idea. More elaborate hybrids use phased execution or hybrid STM, switching between hardware and software modes as contention and capacity behavior change. The important correctness rule is unchanged: the hardware path must synchronize with the fallback path, usually by reading the fallback lock inside the transaction. This is called lock subscription.

The glibc lock-elision history is the warning sign. Eliding locks transparently is attractive, but the best-effort contract, platform errata, and subtle subscription rules make correctness and long-term maintenance difficult. When lock elision is used, it must be treated as a performance path over an ordinary correct synchronization protocol, not as a replacement for that protocol.

Two further interface details complete the picture. \texttt{XTEST} is the introspection instruction of the x86 TSX family: it returns true when executed inside a transactional region (RTM or HLE) and false otherwise, so that \emph{shared} code --- a lock acquisition, a library call --- can detect that it is running speculatively and adjust, most commonly by making lock operations no-ops inside an enclosing transaction instead of aborting it. Any design that routes hardware and software executions through one function body needs exactly this predicate; the companion's runtimes carry the same idea in software form (\texttt{current\_tx->active}). Second, RTM had a sibling interface: Hardware Lock Elision (HLE) marked ordinary \texttt{LOCK}-prefixed instructions with \texttt{xacquire}/\texttt{xrelease} prefixes and asked the processor to elide the lock speculatively, falling back to the plain locked operation on abort. HLE required no new control flow, which made it attractive to libc implementers --- and precisely because the elision was invisible at the source level, programmers could not add subscription reads, retry budgets, or abort accounting to the elided region. That asymmetry, more than any single erratum, is why HLE aged out of the ecosystem while explicit RTM loops survived in HTM-aware codebases.

The availability reality is the final fact, and it motivates this chapter's simulation section. TSX arrived on Haswell (2013) as RTM plus HLE; early steppings had it disabled by microcode once, and in 2021 Intel again shipped microcode disabling TSX on Haswell, Broadwell, and Skylake-family client parts. Cascade Lake server parts shipped with TSX present but disabled by default (enableable in firmware), while recent client generations removed it outright. AMD has never implemented TSX. Arm's TME was specified but never shipped in a commercial core and was withdrawn from the architecture reference manual in 2025~\cite{arm25tmew} --- removed from LLVM~22, the very toolchain Part~V's plugin builds with. Only IBM POWER continues to expose a coherence-based HTM with a mature interface (\texttt{tbegin.}/\texttt{tend.}/\texttt{tabort.}/\texttt{tsr.}, the TEXASR failure-record register, and rollback-only-transaction state). The practical consequence for this book: on most machines you can buy today, none of Table~\ref{tab:htm-implementations}'s first two rows can be executed at all, which is exactly why \S\ref{sec:htm-simulation} exists --- a simulated ISA keeps the instruction set available regardless of what the host implements, and the companion's RDTSC profiling patch (\texttt{patches/profile/tsx}) preserves measurements from the last widely available RTM parts.

\begin{table}[htbp]
\centering
\begin{tabular}{p{0.18\textwidth}p{0.25\textwidth}p{0.25\textwidth}p{0.22\textwidth}}
\hline
\textbf{Mechanism} & \textbf{Tracking} & \textbf{Abort sources} & \textbf{Software obligation}\\
\hline
Intel RTM & Cache-resident speculative read/write sets; explicit \texttt{\_xbegin}, \texttt{\_xend}, \texttt{\_xabort}. & Conflict, capacity, unsupported instruction, interrupt, page fault, explicit abort. & Always provide a correct fallback; subscribe to fallback lock inside the transaction.\\
ARM TME & Architectural transaction begin/end (\texttt{TSTART}/\allowbreak\texttt{TCOMMIT}/\allowbreak\texttt{TCANCEL}/\allowbreak\texttt{TTEST}) with implementation-defined tracking. & Similar best-effort causes; exact diagnostics are implementation-specific. & Same best-effort rule; never assume eventual commit. \emph{TME never shipped in a commercial core and was withdrawn from the Arm ARM in 2025~\cite{arm25tmew}; it is studied here as a design point, not a shipping ISA.}\\
IBM POWER HTM & Coherence-based tracking with a mature architectural interface (\texttt{tbegin.}/\allowbreak\texttt{tend.}/\allowbreak\texttt{tabort.}/\allowbreak\texttt{tsr.}, TEXASR failure recording, ROT). & Conflict, capacity, non-transactional events, implementation limits. & Use retry budgets and fallback paths; exploit richer platform diagnostics where available.\\
gem5 Ruby/HTM model & Simulated cache hierarchy and coherence protocol (\texttt{MESI\_Three\_Level\_HTM}); ISA-agnostic HTM sequencer with calibrated begin/commit latencies. & Commit path is modeled faithfully; conflict/capacity abort wiring is ISA-dependent and incomplete for x86 TSX (see \S\ref{sec:htm-simulation}). & Use simulation to separate effects that hardware counters collapse; check the companion fork's gap list before interpreting abort stats.\\
\hline
\end{tabular}
\caption{HTM implementations differ in interface and diagnostics, but share the same semantic contract: hardware commit is opportunistic and software fallback is mandatory.}
\label{tab:htm-implementations}
\end{table}

\section{Worked Example: RTM, From the Inside}
\index{HTM}\index{RTM}\index{xbegin}\index{xend}
\label{sec:rtm-example}

Intel RTM is the closest commodity hardware comes to a transaction:

\begin{lstlisting}[language=C, caption={The RTM contract: begin, body, fall back on abort.}]
unsigned status = _xbegin();
if (status == _XBEGIN_STARTED) {
    /* transactional region */
    *a -= n; *b += n;            /* both become visible atomically */
    _xend();                     /* commit                        */
} else {
    /* abort: EAX status word encodes the reason (conflict, ...) */
    lock(&g_fallback);           /* best-effort -> serial fallback  */
    *a -= n; *b += n;
    unlock(&g_fallback);
}
\end{lstlisting}

Compare with \S\ref{sec:api-example}: where the STM used a retry loop, RTM uses a hardware fallback instead of a software retry. The \texttt{status} value is a 32-bit word written to \texttt{EAX}: bit~0 set means the abort came from \texttt{\_xabort} (bits 31:24 then carry the \texttt{\_xabort} argument), bit~1 means the transaction \emph{may} succeed on retry, bit~2 a data conflict, bit~3 a capacity overflow, bit~4 a debug breakpoint, bit~5 nesting overflow. It is the entire diagnostics interface --- far poorer than the abort-accounting of Part~VI, which is why profiling tools must infer RTM reasons (Part~VII). Exercise~1 asks for the fallback design; the phased-hybrid and HSTM discussion above is exactly this fallback made adaptive. The transaction's atomicity is enforced by the L1/L2 coherence tracking of the first section --- which is why capacity aborts are unavoidable (see Part~VII's GPU contrast). \textbf{Abort taxonomy}: RTM collapses many causes into the single EAX status word. The practical breakdown is \emph{conflict} (another core touched the read/write set), \emph{capacity} (the transaction outgrew the L1 tracking --- unavoidable at large read/write sets), \emph{unsupported instruction} (e.g.\ most syscalls and \texttt{rdfsbase}), and \emph{debug/explicit} (\texttt{\_xabort} or a debugger single-step). Profilers cannot distinguish conflict from capacity reliably from the status alone --- this is the reason the Part~VII simulator model re-measures them. \textbf{Fallback correctness is the real contract.} The \texttt{lock(\&g\_fallback)} path above is a second, serial execution of the same body. It must be \emph{called with the same state view} as the aborted TSX attempt and must itself be transactional-safe: a crash between the fallback lock and its unlock is indistinguishable from a crash inside the transaction. The failure mode to design against is the \emph{retry loop deadlock}: a hybrid that falls back to a lock but is also invoked by the STM-style retry loop can deadlock on its own fallback mutex (a real bug class in the companion repo's SPHT backend).

The listing is intentionally small, but it hides two non-negotiable rules. First, the fallback body must be the same logical operation as the transactional body. Second, the fallback lock must be visible to the transaction. Without that visibility, the hardware path and fallback path are merely two unsynchronized implementations of the same update.

\section{Worked Example: RTM Bank Transfer with Lock Subscription}
\index{lock subscription}\index{fallback lock}\index{money conservation}
\label{sec:rtm-bank-transfer}

The fallback lock must be visible to the hardware transaction. Otherwise a serialized fallback transfer can race with a hardware transaction that never read the lock.

\begin{lstlisting}[language=C, caption={A compilable RTM bank transfer with account floors and fallback lock subscription.}]
/* Compile on x86-64 with:
 *   c++ -std=c++17 -O2 -mrtm bank_rtm.cc -o bank_rtm
 */
#include <immintrin.h>
#include <atomic>
#include <array>
#include <cassert>
#include <thread>

struct Account {
    long balance;
    long floor;
};

static std::atomic<bool> g_fallback{false};

static void lock_fallback() {
    bool expected = false;
    while (!g_fallback.compare_exchange_weak(
               expected, true,
               std::memory_order_acquire,
               std::memory_order_relaxed)) {
        expected = false;
        std::this_thread::yield();
    }
}

static void unlock_fallback() {
    g_fallback.store(false, std::memory_order_release);
}

static bool transfer_body(Account& from, Account& to, long amount) {
    if (amount < 0) return false;

    __int128 after = static_cast<__int128>(from.balance) - amount;
    if (after < from.floor) return false;

    from.balance -= amount;
    to.balance += amount;
    return true;
}

static bool transfer(Account& from, Account& to, long amount) {
    for (int attempt = 0; attempt < 8; ++attempt) {
        unsigned status = _xbegin();

        if (status == _XBEGIN_STARTED) {
            /*
             * Lock subscription: if the fallback path is active, abort.
             * The load places g_fallback in the transactional read set.
             */
            if (g_fallback.load(std::memory_order_acquire))
                _xabort(0xF0);

            bool ok = transfer_body(from, to, amount);
            _xend();
            return ok;
        }

        if ((status & _XABORT_EXPLICIT) &&
            _XABORT_CODE(status) == 0xF0) {
            std::this_thread::yield();
            continue;
        }

        if ((status & _XABORT_RETRY) == 0)
            break;
    }

    lock_fallback();
    bool ok = transfer_body(from, to, amount);
    unlock_fallback();
    return ok;
}

static long total(const std::array<Account, 2>& a) {
    return a[0].balance + a[1].balance;
}

int main() {
    std::array<Account, 2> bank = {{{100, 0}, {100, 0}}};

    long before = total(bank);
    bool ok1 = transfer(bank[0], bank[1], 25);
    bool ok2 = transfer(bank[0], bank[1], 1000);

    assert(ok1);
    assert(!ok2);
    assert(bank[0].balance >= bank[0].floor);
    assert(bank[1].balance >= bank[1].floor);
    assert(total(bank) == before);
}
\end{lstlisting}

The invariant is money conservation: every successful transfer subtracts and adds the same amount. The account-floor check is inside both the RTM path and the fallback path; fallback is not a different algorithm. The transaction reads \texttt{g\_fallback}; that read is not an optimization but the synchronization point between hardware and software execution. The retry budget is deliberately finite, because a best-effort HTM transaction is allowed to abort forever. Finally, the explicit abort code \texttt{0xF0} separates ``fallback lock is held'' from ordinary conflicts in local diagnostics, though it does not make RTM deterministic.

The companion backends use this same pattern at different levels. SGL serializes everything. TSX-SGL attempts the critical section in hardware and falls back to serialization. Software systems such as TL2, NOrec, TSC-TM, SwissTM, and TinySTM do not depend on cache-resident speculation, but they pay for metadata, validation, and logging. The measured results in the companion project reflect that trade-off: for the fuzz/bank workload on an M1 Pro, TL2 and TSC-TM are around two million transactions per second, while NOrec-BF improves the read-mostly NOrec path by reducing validation cost. These numbers are workload-specific; they are useful here only as a reminder that HTM is one point in a larger design space.

\section{Worked Example: Abort-Cause Counters}
\label{sec:htm-abort-counters}

RTM exposes limited status bits. A useful first diagnostic is therefore local accounting: count commits, explicit subscription aborts, retryable aborts, and non-retryable aborts. The counters must be updated outside the transaction. If the transaction aborts, writes performed inside it vanish; if the writes force extra cache lines into the transaction, they also perturb the abort rate being measured.

\begin{lstlisting}[language=C, caption={Per-thread RTM counters kept outside the transactional footprint.}]
/* Compile on x86-64 with:
 *   c++ -std=c++17 -O2 -mrtm rtm_counters.cc -o rtm_counters
 */
#include <immintrin.h>
#include <atomic>
#include <cassert>
#include <cstdint>

struct Counters {
    std::uint64_t commits = 0;
    std::uint64_t explicit_lock = 0;
    std::uint64_t retryable = 0;
    std::uint64_t nonretryable = 0;
};

static std::atomic<bool> g_lock{false};

static void critical_increment(int& x) {
    ++x;
}

static void run_once(int& x, Counters& c) {
    for (int i = 0; i != 4; ++i) {
        unsigned s = _xbegin();
        if (s == _XBEGIN_STARTED) {
            if (g_lock.load(std::memory_order_acquire))
                _xabort(0xF0);
            critical_increment(x);
            _xend();
            ++c.commits;
            return;
        }

        if ((s & _XABORT_EXPLICIT) && _XABORT_CODE(s) == 0xF0) {
            ++c.explicit_lock;
        } else if (s & _XABORT_RETRY) {
            ++c.retryable;
        } else {
            ++c.nonretryable;
            break;
        }
    }

    while (g_lock.exchange(true, std::memory_order_acquire)) {
    }
    critical_increment(x);
    g_lock.store(false, std::memory_order_release);
}

int main() {
    int x = 0;
    Counters c;
    run_once(x, c);
    assert(x == 1);
    assert(c.commits + c.explicit_lock + c.retryable + c.nonretryable >= 1);
}
\end{lstlisting}

Counters are diagnostic state, not transactional state: a counter increment inside RTM would be rolled back on abort and would add another write line on commit. The status bits classify symptoms; they do not fully explain the microarchitectural cause.

\section{Worked Example: Capacity Probe}
\label{sec:htm-capacity-probe}

The following program sweeps the number of cache lines touched by a transaction. It is not a benchmark result; it is a harness for observing the qualitative transition from small, cache-resident transactions to transactions that frequently abort for capacity or other implementation limits.

\begin{lstlisting}[language=C, caption={A small RTM capacity probe that varies the transactional footprint.}]
/* Compile on x86-64 with:
 *   c++ -std=c++17 -O2 -mrtm rtm_capacity.cc -o rtm_capacity
 */
#include <immintrin.h>
#include <array>
#include <cassert>
#include <cstddef>
#include <iostream>

struct alignas(64) Line {
    int value;
    char pad[60];
};

int main() {
    std::array<Line, 256> lines{};
    for (std::size_t touched = 1; touched <= lines.size(); touched *= 2) {
        unsigned status = _xbegin();
        if (status == _XBEGIN_STARTED) {
            for (std::size_t i = 0; i < touched; ++i)
                lines[i].value++;
            _xend();
            std::cout << touched << " commit\n";
        } else {
            std::cout << touched << " abort status=" << status << "\n";
        }
    }

    long sum = 0;
    for (const auto& l : lines)
        sum += l.value;
    assert(sum >= 0);
}
\end{lstlisting}

Aligning each element to a cache line makes the footprint easier to reason about. The program does not assume a particular capacity threshold. The same idea is useful in simulation, where cache sizes and associativity can be varied deliberately.

\section{Simulation with gem5}
\label{sec:htm-simulation}

Simulation answers questions that commodity counters often merge together. A hardware counter may report an abort, or the EAX status word may hint that the abort was retryable, but it rarely tells whether the root cause was an L1 conflict, an eviction, a coherence race, a TLB event, or an implementation limit. A simulator can expose the cache hierarchy, coherence protocol, and transactional metadata directly.

Simulation has a second role for this book's audience: \emph{running TM code on machines that do not implement a TM instruction set at all.} TSX is unavailable on AMD CPUs and has been removed or disabled on many recent Intel parts; TME never shipped and was withdrawn from the Arm ARM in 2025 (it is gone from LLVM~22, the very toolchain used by Part~V's plugin). A simulated ISA makes the instruction set available regardless of the host.

The companion gem5 fork (the \textsc{TM-SIM} repository, gem5 v25.1~\cite{binkert11gem5} plus an in-tree x86 TSX implementation) is the concrete vehicle. Upstream gem5 contains \emph{no} x86 TSX support; the fork adds the decoder and instruction semantics for \texttt{XBEGIN}/\texttt{XEND}/\texttt{XABORT}/\texttt{XTEST} and drives them through Ruby's \texttt{MESI\_Three\_Level\_HTM} protocol, whose ISA-agnostic \texttt{HTMSequencer} is calibrated with the measured Broadwell-EP costs of Chapter~14's machine profiles (\texttt{htm\_start\_\allowbreak latency}\,\allowbreak${}=60$ and \texttt{htm\_commit\_\allowbreak latency}\,\allowbreak${}=178$ cycles, matching the RDTSC-profiled XBEGIN/XEND costs from real hardware). The workflow is:

\begin{enumerate}
    \item build the TSX target: \texttt{scons build/X86\_TSX/gem5.opt --with-ruby};
    \item cross-compile a static benchmark against a TSX-consuming backend (\texttt{make BACKEND=TSXSGL GEM5=1} in the companion repository --- the host needs no TSX; only the simulated CPU executes it);
    \item run in syscall-emulation mode with a bounded, ROI-delimited workload (\texttt{bank -n 500}), with \texttt{m5} statistics reset/dump markers bracketing the transactional phase;
    \item read per-transaction cycles and read/write-set histograms from the Ruby HTM sequencer statistics;
    \item for full-system fidelity, boot once, checkpoint, and restore with the detailed CPU model so that only the region of interest is simulated slowly.
\end{enumerate}

\paragraph{Current limits of the fork (honest status).} The calibration target met so far is single-threaded commit-path fidelity: on the bank microbenchmark the simulated per-transaction cost lands within $\sim$1\% of the cycle cost model. Three gaps are documented in the fork's \texttt{IMPROVEMENT\_PLAN.md}: \texttt{XABORT} currently raises \texttt{InvalidOpcode} instead of delivering the EAX status; conflict and capacity aborts are not yet wired from the L0 cache's \texttt{LD\_FAIL}/\texttt{ST\_FAIL} callbacks to the XBEGIN fallback; and the abort path has only been exercised single-threaded. Until those land, multi-threaded abort behavior should be studied through the repository's real-hardware RDTSC profiling patch (\texttt{patches/profile/tsx}) and the Rust \texttt{tsx\_sim} cost-model backend, and gem5 results should be read as commit-path numbers.

The capacity probe above becomes more informative in simulation than on a fixed machine --- once the capacity-abort wiring lands. On real hardware, the abort boundary is affected by undocumented structures and unrelated activity. In simulation, the experimenter can hold the workload constant and vary L1 size, associativity, line size, or contention. This makes transient effects visible: evictions mid-transaction, write-set coalescing, false conflicts inside a coherence domain, and instruction events that force an abort.

The simulator is not a replacement for measurement. It is a microscope. It explains mechanisms under controlled assumptions; the hardware run tells whether those mechanisms dominate on a real platform. The three instruments used together --- real-hardware profiling patches, the gem5 fork, and the Rust cost-model simulator --- cover, respectively, ground truth, mechanism, and speed.

\section{Worked Example: Model Checking HTM Abort and Fallback}
\index{TLA+}\index{abort}\index{fallback path}
\label{sec:htm-fallback-tla}

A small model is enough to expose the essential rule: abort may occur at any time before commit, and fallback must preserve the same invariants as the hardware path.

\begin{lstlisting}[language={}, caption={TLA+ model of one transferring transaction, abort, interference, and fallback.}]
-------------------------- MODULE BankHTMFallback --------------------------
EXTENDS Integers, TLC

Accounts == {0, 1}
InitA == 100
InitB == 100
Floor == 0
Amt == 7

VARIABLES bal, pc, snap

vars == <<bal, pc, snap>>

Total(b) == b[0] + b[1]

CanMove(b) == b[0] - Amt >= Floor

Move(b) == [b EXCEPT ![0] = @ - Amt,
                      ![1] = @ + Amt]

Init ==
    /\ bal = [a \in Accounts |-> IF a = 0 THEN InitA ELSE InitB]
    /\ pc = "idle"
    /\ snap = bal

Begin ==
    /\ pc = "idle"
    /\ pc' = "tx"
    /\ snap' = bal
    /\ UNCHANGED bal

Abort ==
    /\ pc = "tx"
    /\ pc' = "idle"
    /\ UNCHANGED <<bal, snap>>

Commit ==
    /\ pc = "tx"
    /\ bal = snap
    /\ CanMove(snap)
    /\ bal' = Move(bal)
    /\ pc' = "idle"
    /\ snap' = snap

CommitNoFunds ==
    /\ pc = "tx"
    /\ bal = snap
    /\ ~CanMove(snap)
    /\ pc' = "idle"
    /\ UNCHANGED <<bal, snap>>

Fallback ==
    /\ pc = "idle"
    /\ CanMove(bal)
    /\ bal' = Move(bal)
    /\ pc' = "idle"
    /\ UNCHANGED snap

InterfereAbort ==
    /\ pc = "tx"
    /\ CanMove(bal)
    /\ bal' = Move(bal)
    /\ pc' = "idle"
    /\ snap' = snap

Next ==
    \/ Begin
    \/ Abort
    \/ Commit
    \/ CommitNoFunds
    \/ Fallback
    \/ InterfereAbort

Spec == Init /\ [][Next]_vars

MoneyConserved == Total(bal) = InitA + InitB

FloorSafe == \A a \in Accounts : bal[a] >= Floor

THEOREM Spec => []MoneyConserved
THEOREM Spec => []FloorSafe
=============================================================================
\end{lstlisting}

\texttt{Commit} requires \texttt{bal = snap}; this is the abstract coherence check --- if the snapshot is stale, hardware must not commit it. \texttt{InterfereAbort} models another successful transfer that conflicts with the running transaction; the running transaction aborts, but money is still conserved. \texttt{Fallback} executes only from \texttt{idle}, so it is serialized with respect to the modeled hardware transaction. Both invariants are safety properties; they do not claim that a best-effort transaction eventually commits.

The model omits caches, lines, and invalidations. That is intentional. At the proof boundary, HTM contributes one rule: a transaction may commit only if its snapshot is still valid. Everything else is fallback correctness. This is why the companion project model-checks the software contract as well as testing the implementations.

\section{Summary}

HTM leverages existing speculation and cache coherence to make short transactions fast. A commit can be nearly as cheap as leaving a speculative region, and conflict detection can be piggybacked on invalidation traffic already present in the memory hierarchy.

The cost is uncertainty. Transactions can abort for conflicts, capacity, interrupts, faults, unsupported instructions, and implementation limits. Therefore every HTM design is also a fallback design. Correct programs subscribe to the fallback lock, bound their retries, keep diagnostics outside the transactional footprint, and validate the same invariants on the hardware and software paths.

\section{Exercises}
\begin{enumerate}
    \item \emph{[implementation]} In Intel RTM, a transaction that exceeds the L1 cache capacity always aborts. Propose a software fallback path and discuss its performance implications. If you do not have RTM hardware, ground the design in the companion's TSX-SGL backend (\texttt{backends/tm\_impl/tsx\_sgl}) and the phased-hybrid model of \S\ref{sec:hybrid-example}.
    \item \emph{[theory]} Using a conceptual pipeline diagram, show why a transaction that triggers a TLB miss is likely to abort.
    \item \emph{[implementation]} Footprint sweep: run the capacity probe of \S\ref{sec:htm-capacity-probe} on RTM hardware with the companion's RDTSC profiling patch (\texttt{patches/profile/tsx}), or on any host against the companion's TSX simulation backend (\texttt{simulator/}, \texttt{runtime/tsx\_sim}), and plot the abort rate as the number of touched cache lines grows. Explain the trend. (The gem5 fork of \S\ref{sec:htm-simulation} cannot yet answer this one: its capacity-abort wiring is an open gap.)
    \item \emph{[implementation]} Extend Listing~\ref{sec:rtm-bank-transfer}'s transfer code to collect per-thread counters for commits, explicit fallback-lock aborts, retryable aborts, and non-retryable aborts. Keep the counters outside the transaction and explain why updating them inside RTM would distort the measured abort rate.
    \item \emph{[verification]} Modify the TLA+ model in \S\ref{sec:htm-fallback-tla} so that transfers may go in both directions. State and check the money-conservation invariant and an account-floor invariant for both accounts.
    \item \emph{[theory]} Explain why lock subscription is necessary for opacity in a hybrid HTM plus lock fallback design. Give an execution in which omitting the subscription lets a hardware transaction commit concurrently with the fallback critical section.
    \item \emph{[implementation]} Modify the capacity probe so that it first reads many cache lines and then writes only two. Compare the expected abort behavior with a write-heavy transaction that writes every touched line.
    \item \emph{[verification]} Add a boolean fallback-lock variable to the TLA+ model. Permit \texttt{Begin} only when the lock is clear, and add an explicit abort transition when the lock becomes held during a transaction.
    \item \emph{[theory]} Explain why a best-effort HTM interface cannot provide wait-freedom by itself, even when every individual transaction is finite and conflict-free in the source program.
    \item \emph{[literature]} Read the original Herlihy--Moss proposal
          \cite{herlihy93} and Lev, Moir, and Nussbaum's
          \emph{Phased Transactional Memory} \cite{lev07phased}.  Explain
          how the phased hybrid generalizes the raw XCOMMIT/XABORT design:
          which part of Herlihy--Moss survives, which part is replaced by
          a software fallback, and what the phase-change criterion buys
          over a fixed subscription lock.
    \item \emph{[open]} Best-effort HTM moved the fallback obligation into
          software and that decision let TSX ship.  Argue a position:
          should the field push hardware toward \emph{stronger} guarantees
          (progress, capacity, deterministic aborts), or is best-effort
          the right abstraction and the remaining work purely in software?
          Tie your position to at least one mechanism from this chapter
          (capacity, conflict, fallback lock).
\end{enumerate}

\textbf{How to check your work.} The \emph{[implementation]} exercises
verify on the companion: the fallback design of Exercise~1 maps onto
\texttt{backends/tm\_impl/tsx\_sgl} (TSX path with lock subscription);
Exercise~4 is checked with the abort counters of
\S\ref{sec:htm-abort-counters}; and Exercise~3/7 run the capacity probe of
\S\ref{sec:htm-capacity-probe} on real RTM hardware or against the
\texttt{tsx\_sim} cost-model backend.  The \emph{[verification]} exercises are
checked with TLC against the model of \S\ref{sec:htm-fallback-tla} using
the workflow of Appendix~\ref{chap:pluscal-quick-ref}.  The
\emph{[theory]} exercises are graded against the definition of
\S\ref{sec:rtm-example} and the subscription argument of
\S\ref{sec:htm-fallback-tla}.  The \emph{[literature]} exercise is graded
on whether your ``survives / replaced'' split matches the two papers.  The
\emph{[open]} exercise has no single right answer; it is graded on whether
your position is tied to a concrete mechanism from this chapter.

\index{Herlihy and Moss}
\index{XCOMMIT}
\index{XABORT}
\index{MESI}
\index{speculative stores}
\index{transactional read set}
\index{transactional write set}
\index{coherence invalidation}
\index{capacity abort}
\index{conflict abort}
\index{unsupported instruction abort}
\index{Intel TSX}
\index{Intel RTM}
\index{ARM TME}
\index{IBM POWER HTM}
\index{XTEST}
\index{HLE}
\index{xacquire}
\index{xrelease}
\index{TSX availability}
\index{gem5}
\index{Ruby memory system}
\index{lock elision}
\index{TSX-SGL}
\index{phased hybrid TM}
\index{hybrid STM}
\index{abort counters}
\index{cache footprint}
\index{fallback serialization}
\index{opacity!HTM fallback}
\index{bank transfer invariant}
% -------------------------------------------------------------------
